\chapter*{Résumé}\label{chap:00-resume}
\addcontentsline{toc}{chapter}{Résumé}

\paragraph{La question}
En 2018, à ICLR puis à NeurIPS, Joelle Pineau a montré pourquoi les résultats publiés en apprentissage
par renforcement profond (RL) inspiraient peu confiance. Nous y relevons seize failles et demandons, pour chacune, si elle a été écoutée.

\paragraph{La méthode}
Deux instruments servent, sous des protocoles écrits d'avance dont chaque écart est daté. Le
sondage relit des articles avec la grille de Pineau (Q1 à Q8), sept critères tirés de sa conférence
d'ICLR (X1 à X7) et une seconde grille (Y1 à Y12). Des modèles d'IA agentique lisent chaque article
entier et justifient chaque décision par une citation paginée. Le verdict compare 50 articles de RL
classique de 2026 (hors modèles de langue) à notre lecture de 50 articles de 2018, avec les mêmes
grilles et les mêmes lecteurs. Vingt articles par année lue entre 2019 et 2024 datent le changement. Les expériences refont celles de 2018
avec PPO (en PyTorch et en JAX), SAC et CrossQ sur quatre tâches MuJoCo, dix graines aléatoires par
programme et par tâche, puis vingt. Un volet entraîne deux petits modèles de langue avec GRPO.

\paragraph{Les résultats du sondage}
Notre lecture de 2018 s'écarte des chiffres de Pineau : 55~\% des articles donnent tous les
hyperparamètres de leur méthode, quand elle en comptait 90~\%. Les verdicts comparent donc nos deux
lectures. Dans le RL classique, quatre pratiques ont changé : le lien vers le code
passe de 31~\% à 74~\% ($+43$ points, $[+23\,;\,+58]$), la sensibilité à un réglage montrée de 44~\% à
80~\%, le calcul chiffré de 27~\% à 58~\%, une incertitude rapportée de 63~\% à 88~\%. Cette dernière
est fragile : sans treize articles lus en partie par un autre modèle de langue, elle échoue à la
correction de Holm ($p_{\text{Holm}} = 0{,}07$). Les quatre montent lentement, sans saut après la liste de
contrôle de NeurIPS de 2019. Le test statistique (de 6~\% à 12~\%), les références décrites
avec tous leurs réglages (35~\% contre 29~\%), l'évaluation tenue à part (33~\% puis 45~\%) et la
version des bancs d'essai (23~\% puis 26~\%) n'ont pas bougé de façon détectable. Pineau a été écoutée
sur ce qui se publie à côté du résultat, beaucoup moins sur la manière de le juger.

\paragraph{Les résultats des expériences}
Sur trente essais de PPO, deux groupes de cinq tirés au hasard trompent un
test de Welch dans 1,2~\% des cas, triés dans 31~\%. Garder les trois meilleurs essais sur dix gonfle
la moyenne de 33~\%. Le script de Peter Henderson, qui ignore la dispersion entre graines aléatoires,
explique très probablement le $t$ de 9 de 2018 : sur nos essais de TRPO, il dépasse 9,09 dans 38~\% des tirages. À budget égal, avec
vingt graines aléatoires, CrossQ mène sur les quatre tâches. Une physique modifiée de 20~\%
envoie le meneur au dernier rang sur \task{Hopper}. Sur Qwen, avec cinq graines
aléatoires, GRPO gagne $+14$ points, intervalle de $-3$ à $+34$, qui contient zéro.

\paragraph{Code et données}
Dépôt \repo{} (BSD 3-Clause) : protocoles et lectures citées (\file{study/}), entraînements (\file{sim/}),
statistiques (\file{src/pineau\_rl/}). \file{make verify} régénère et vérifie chaque résultat.
